% sections/introduction.tex
% Included from main.tex via % sections/introduction.tex
% Included from main.tex via % sections/introduction.tex
% Included from main.tex via % sections/introduction.tex
% Included from main.tex via \input{sections/introduction}

\section{Introduction}
\label{sec:introduction}

Three converging frontiers are reshaping autonomous systems research.
The \emph{low-altitude economy} demands scalable infrastructure for urban air
mobility, drone logistics, and aerial inspection.
\emph{Embodied intelligence} requires agents that perceive and act in shared
physical environments through vision, language, and control.
\emph{Air-ground cooperative systems} bring these threads together, calling for
heterogeneous robots that operate jointly across aerial and ground domains.
Simulation is essential for advancing all three frontiers, as real-world
deployment is costly, safety-critical, and difficult to scale.
Yet no widely adopted open-source platform provides a unified infrastructure
capable of jointly modeling aerial and ground agents within a single physically
coherent environment.
CARLA-Air is designed to fill this gap.

\paragraph{The simulation landscape and its gap.}
Existing open-source simulators address complementary domains without overlap.
CARLA~\cite{carla}, built on Unreal Engine~4~\cite{ue4}, has become the de
facto standard for urban autonomous driving research, offering photorealistic
environments, rich traffic populations, and a mature Python API.
AirSim~\cite{airsim}, also built on UE4, provides physics-accurate multirotor
simulation with high-frequency dynamics and a comprehensive aerial sensor suite.
The limitation of each platform is precisely the strength of the other: CARLA
lacks aerial agents, while AirSim lacks realistic ground traffic and pedestrian
interactions.
Meanwhile, AirSim's upstream development has been archived by its original
maintainers, leaving a widely adopted flight simulation stack without an active
evolution path.
Other simulators across driving, UAV, and embodied AI domains similarly remain
confined to a single agent modality (see Section~\ref{sec:related} for a
comprehensive survey).
As a result, emerging workloads that span both air and ground
domains---air-ground cooperation, cross-domain embodied navigation, joint
multi-modal data collection, and cooperative reinforcement learning---lack a
shared simulation foundation.

\paragraph{Why not bridge-based co-simulation?}
A common workaround connects heterogeneous simulators through bridge-based
co-simulation, typically via ROS\,2~\cite{ros2} or custom message-passing
interfaces.
While functionally viable, such approaches introduce inter-process
synchronization complexity, communication overhead, and duplicated rendering
pipelines.
More critically, independent simulation processes cannot guarantee strict
spatial-temporal consistency across sensor streams---a requirement for
perception, learning, and evaluation in embodied intelligence systems.
Fig.~\ref{fig:ipc_overhead} quantifies the per-frame inter-process overhead
contrast between bridge-based co-simulation and the single-process design
adopted by CARLA-Air.

\input{figures/ipc_plot}

\paragraph{CARLA-Air: a unified infrastructure for air-ground embodied intelligence.}
We present \textbf{CARLA-Air}, an open-source platform that integrates CARLA
and AirSim within a single Unreal Engine process, purpose-built as a practical
simulation foundation for air-ground embodied intelligence research.
By inheriting and extending AirSim's aerial simulation capabilities within a
modern, actively maintained infrastructure, CARLA-Air also provides a
sustainable evolution path for the large body of existing AirSim-based research.
Key capabilities of the platform include:

\begin{enumerate}[leftmargin=*, label=(\roman*)]

\item \textbf{Single-process air-ground integration.}
CARLA-Air resolves a fundamental engine-level conflict---UE4 permits only one
active game mode per world---through a composition-based design that inherits
all ground simulation subsystems from CARLA while spawning AirSim's aerial
flight actor as a regular world entity.
This yields a shared physics tick, a shared rendering pipeline, and strict
spatial-temporal consistency across all sensor viewpoints.

\item \textbf{Full API compatibility and zero-modification code migration.}
Both CARLA and AirSim native Python APIs and ROS\,2 interfaces are fully
preserved, allowing existing research codebases to run on CARLA-Air without
modification.

\item \textbf{Photorealistic, physically coherent simulation world.}
The platform delivers rich urban and natural environments populated with
rule-compliant traffic flow, socially-aware pedestrians, and aerodynamically
consistent multirotor dynamics, with synchronized capture of up to 18 sensor
modalities across all aerial and ground platforms at each simulation tick.

\item \textbf{Extensible asset pipeline.}
Researchers can import custom robot platforms, UAV configurations, vehicles,
and environment maps into the shared simulation world, enabling flexible
construction of diverse air-ground interaction scenarios.

\end{enumerate}

\noindent Building on these capabilities, CARLA-Air provides out-of-the-box
support for representative air-ground embodied intelligence workloads across
four research directions:

\begin{enumerate}[leftmargin=*, label=(\alph*)]

\item \textbf{Air-ground cooperation}---heterogeneous aerial and ground agents
coordinate within a shared environment for tasks such as cooperative
surveillance, escort, and search-and-rescue.

\item \textbf{Embodied navigation and vision-language action}---agents navigate
and act grounded in visual and linguistic input, leveraging both aerial
overview and ground-level detail.

\item \textbf{Multi-modal perception and dataset construction}---synchronized
aerial-ground sensor streams are collected at scale to build paired datasets
for cross-view perception, 3D reconstruction, and scene understanding.

\item \textbf{Reinforcement-learning-based policy training}---agents learn
cooperative or individual policies through closed-loop interaction in
physically consistent air-ground environments.

\end{enumerate}

\noindent As a lightweight and practical infrastructure, CARLA-Air lowers the
barrier for developing and evaluating air-ground embodied intelligence systems,
and provides a unified simulation foundation for emerging applications in
low-altitude robotics, cross-domain autonomy, and large-scale embodied AI
research.



\section{Introduction}
\label{sec:introduction}

Three converging frontiers are reshaping autonomous systems research.
The \emph{low-altitude economy} demands scalable infrastructure for urban air
mobility, drone logistics, and aerial inspection.
\emph{Embodied intelligence} requires agents that perceive and act in shared
physical environments through vision, language, and control.
\emph{Air-ground cooperative systems} bring these threads together, calling for
heterogeneous robots that operate jointly across aerial and ground domains.
Simulation is essential for advancing all three frontiers, as real-world
deployment is costly, safety-critical, and difficult to scale.
Yet no widely adopted open-source platform provides a unified infrastructure
capable of jointly modeling aerial and ground agents within a single physically
coherent environment.
CARLA-Air is designed to fill this gap.

\paragraph{The simulation landscape and its gap.}
Existing open-source simulators address complementary domains without overlap.
CARLA~\cite{carla}, built on Unreal Engine~4~\cite{ue4}, has become the de
facto standard for urban autonomous driving research, offering photorealistic
environments, rich traffic populations, and a mature Python API.
AirSim~\cite{airsim}, also built on UE4, provides physics-accurate multirotor
simulation with high-frequency dynamics and a comprehensive aerial sensor suite.
The limitation of each platform is precisely the strength of the other: CARLA
lacks aerial agents, while AirSim lacks realistic ground traffic and pedestrian
interactions.
Meanwhile, AirSim's upstream development has been archived by its original
maintainers, leaving a widely adopted flight simulation stack without an active
evolution path.
Other simulators across driving, UAV, and embodied AI domains similarly remain
confined to a single agent modality (see Section~\ref{sec:related} for a
comprehensive survey).
As a result, emerging workloads that span both air and ground
domains---air-ground cooperation, cross-domain embodied navigation, joint
multi-modal data collection, and cooperative reinforcement learning---lack a
shared simulation foundation.

\paragraph{Why not bridge-based co-simulation?}
A common workaround connects heterogeneous simulators through bridge-based
co-simulation, typically via ROS\,2~\cite{ros2} or custom message-passing
interfaces.
While functionally viable, such approaches introduce inter-process
synchronization complexity, communication overhead, and duplicated rendering
pipelines.
More critically, independent simulation processes cannot guarantee strict
spatial-temporal consistency across sensor streams---a requirement for
perception, learning, and evaluation in embodied intelligence systems.
Fig.~\ref{fig:ipc_overhead} quantifies the per-frame inter-process overhead
contrast between bridge-based co-simulation and the single-process design
adopted by CARLA-Air.

\input{figures/ipc_plot}

\paragraph{CARLA-Air: a unified infrastructure for air-ground embodied intelligence.}
We present \textbf{CARLA-Air}, an open-source platform that integrates CARLA
and AirSim within a single Unreal Engine process, purpose-built as a practical
simulation foundation for air-ground embodied intelligence research.
By inheriting and extending AirSim's aerial simulation capabilities within a
modern, actively maintained infrastructure, CARLA-Air also provides a
sustainable evolution path for the large body of existing AirSim-based research.
Key capabilities of the platform include:

\begin{enumerate}[leftmargin=*, label=(\roman*)]

\item \textbf{Single-process air-ground integration.}
CARLA-Air resolves a fundamental engine-level conflict---UE4 permits only one
active game mode per world---through a composition-based design that inherits
all ground simulation subsystems from CARLA while spawning AirSim's aerial
flight actor as a regular world entity.
This yields a shared physics tick, a shared rendering pipeline, and strict
spatial-temporal consistency across all sensor viewpoints.

\item \textbf{Full API compatibility and zero-modification code migration.}
Both CARLA and AirSim native Python APIs and ROS\,2 interfaces are fully
preserved, allowing existing research codebases to run on CARLA-Air without
modification.

\item \textbf{Photorealistic, physically coherent simulation world.}
The platform delivers rich urban and natural environments populated with
rule-compliant traffic flow, socially-aware pedestrians, and aerodynamically
consistent multirotor dynamics, with synchronized capture of up to 18 sensor
modalities across all aerial and ground platforms at each simulation tick.

\item \textbf{Extensible asset pipeline.}
Researchers can import custom robot platforms, UAV configurations, vehicles,
and environment maps into the shared simulation world, enabling flexible
construction of diverse air-ground interaction scenarios.

\end{enumerate}

\noindent Building on these capabilities, CARLA-Air provides out-of-the-box
support for representative air-ground embodied intelligence workloads across
four research directions:

\begin{enumerate}[leftmargin=*, label=(\alph*)]

\item \textbf{Air-ground cooperation}---heterogeneous aerial and ground agents
coordinate within a shared environment for tasks such as cooperative
surveillance, escort, and search-and-rescue.

\item \textbf{Embodied navigation and vision-language action}---agents navigate
and act grounded in visual and linguistic input, leveraging both aerial
overview and ground-level detail.

\item \textbf{Multi-modal perception and dataset construction}---synchronized
aerial-ground sensor streams are collected at scale to build paired datasets
for cross-view perception, 3D reconstruction, and scene understanding.

\item \textbf{Reinforcement-learning-based policy training}---agents learn
cooperative or individual policies through closed-loop interaction in
physically consistent air-ground environments.

\end{enumerate}

\noindent As a lightweight and practical infrastructure, CARLA-Air lowers the
barrier for developing and evaluating air-ground embodied intelligence systems,
and provides a unified simulation foundation for emerging applications in
low-altitude robotics, cross-domain autonomy, and large-scale embodied AI
research.



\section{Introduction}
\label{sec:introduction}

Three converging frontiers are reshaping autonomous systems research.
The \emph{low-altitude economy} demands scalable infrastructure for urban air
mobility, drone logistics, and aerial inspection.
\emph{Embodied intelligence} requires agents that perceive and act in shared
physical environments through vision, language, and control.
\emph{Air-ground cooperative systems} bring these threads together, calling for
heterogeneous robots that operate jointly across aerial and ground domains.
Simulation is essential for advancing all three frontiers, as real-world
deployment is costly, safety-critical, and difficult to scale.
Yet no widely adopted open-source platform provides a unified infrastructure
capable of jointly modeling aerial and ground agents within a single physically
coherent environment.
CARLA-Air is designed to fill this gap.

\paragraph{The simulation landscape and its gap.}
Existing open-source simulators address complementary domains without overlap.
CARLA~\cite{carla}, built on Unreal Engine~4~\cite{ue4}, has become the de
facto standard for urban autonomous driving research, offering photorealistic
environments, rich traffic populations, and a mature Python API.
AirSim~\cite{airsim}, also built on UE4, provides physics-accurate multirotor
simulation with high-frequency dynamics and a comprehensive aerial sensor suite.
The limitation of each platform is precisely the strength of the other: CARLA
lacks aerial agents, while AirSim lacks realistic ground traffic and pedestrian
interactions.
Meanwhile, AirSim's upstream development has been archived by its original
maintainers, leaving a widely adopted flight simulation stack without an active
evolution path.
Other simulators across driving, UAV, and embodied AI domains similarly remain
confined to a single agent modality (see Section~\ref{sec:related} for a
comprehensive survey).
As a result, emerging workloads that span both air and ground
domains---air-ground cooperation, cross-domain embodied navigation, joint
multi-modal data collection, and cooperative reinforcement learning---lack a
shared simulation foundation.

\paragraph{Why not bridge-based co-simulation?}
A common workaround connects heterogeneous simulators through bridge-based
co-simulation, typically via ROS\,2~\cite{ros2} or custom message-passing
interfaces.
While functionally viable, such approaches introduce inter-process
synchronization complexity, communication overhead, and duplicated rendering
pipelines.
More critically, independent simulation processes cannot guarantee strict
spatial-temporal consistency across sensor streams---a requirement for
perception, learning, and evaluation in embodied intelligence systems.
Fig.~\ref{fig:ipc_overhead} quantifies the per-frame inter-process overhead
contrast between bridge-based co-simulation and the single-process design
adopted by CARLA-Air.

\input{figures/ipc_plot}

\paragraph{CARLA-Air: a unified infrastructure for air-ground embodied intelligence.}
We present \textbf{CARLA-Air}, an open-source platform that integrates CARLA
and AirSim within a single Unreal Engine process, purpose-built as a practical
simulation foundation for air-ground embodied intelligence research.
By inheriting and extending AirSim's aerial simulation capabilities within a
modern, actively maintained infrastructure, CARLA-Air also provides a
sustainable evolution path for the large body of existing AirSim-based research.
Key capabilities of the platform include:

\begin{enumerate}[leftmargin=*, label=(\roman*)]

\item \textbf{Single-process air-ground integration.}
CARLA-Air resolves a fundamental engine-level conflict---UE4 permits only one
active game mode per world---through a composition-based design that inherits
all ground simulation subsystems from CARLA while spawning AirSim's aerial
flight actor as a regular world entity.
This yields a shared physics tick, a shared rendering pipeline, and strict
spatial-temporal consistency across all sensor viewpoints.

\item \textbf{Full API compatibility and zero-modification code migration.}
Both CARLA and AirSim native Python APIs and ROS\,2 interfaces are fully
preserved, allowing existing research codebases to run on CARLA-Air without
modification.

\item \textbf{Photorealistic, physically coherent simulation world.}
The platform delivers rich urban and natural environments populated with
rule-compliant traffic flow, socially-aware pedestrians, and aerodynamically
consistent multirotor dynamics, with synchronized capture of up to 18 sensor
modalities across all aerial and ground platforms at each simulation tick.

\item \textbf{Extensible asset pipeline.}
Researchers can import custom robot platforms, UAV configurations, vehicles,
and environment maps into the shared simulation world, enabling flexible
construction of diverse air-ground interaction scenarios.

\end{enumerate}

\noindent Building on these capabilities, CARLA-Air provides out-of-the-box
support for representative air-ground embodied intelligence workloads across
four research directions:

\begin{enumerate}[leftmargin=*, label=(\alph*)]

\item \textbf{Air-ground cooperation}---heterogeneous aerial and ground agents
coordinate within a shared environment for tasks such as cooperative
surveillance, escort, and search-and-rescue.

\item \textbf{Embodied navigation and vision-language action}---agents navigate
and act grounded in visual and linguistic input, leveraging both aerial
overview and ground-level detail.

\item \textbf{Multi-modal perception and dataset construction}---synchronized
aerial-ground sensor streams are collected at scale to build paired datasets
for cross-view perception, 3D reconstruction, and scene understanding.

\item \textbf{Reinforcement-learning-based policy training}---agents learn
cooperative or individual policies through closed-loop interaction in
physically consistent air-ground environments.

\end{enumerate}

\noindent As a lightweight and practical infrastructure, CARLA-Air lowers the
barrier for developing and evaluating air-ground embodied intelligence systems,
and provides a unified simulation foundation for emerging applications in
low-altitude robotics, cross-domain autonomy, and large-scale embodied AI
research.



\section{Introduction}
\label{sec:introduction}

Three converging frontiers are reshaping autonomous systems research.
The \emph{low-altitude economy} demands scalable infrastructure for urban air
mobility, drone logistics, and aerial inspection.
\emph{Embodied intelligence} requires agents that perceive and act in shared
physical environments through vision, language, and control.
\emph{Air-ground cooperative systems} bring these threads together, calling for
heterogeneous robots that operate jointly across aerial and ground domains.
Simulation is essential for advancing all three frontiers, as real-world
deployment is costly, safety-critical, and difficult to scale.
Yet no widely adopted open-source platform provides a unified infrastructure
capable of jointly modeling aerial and ground agents within a single physically
coherent environment.
CARLA-Air is designed to fill this gap.

\paragraph{The simulation landscape and its gap.}
Existing open-source simulators address complementary domains without overlap.
CARLA~\cite{carla}, built on Unreal Engine~4~\cite{ue4}, has become the de
facto standard for urban autonomous driving research, offering photorealistic
environments, rich traffic populations, and a mature Python API.
AirSim~\cite{airsim}, also built on UE4, provides physics-accurate multirotor
simulation with high-frequency dynamics and a comprehensive aerial sensor suite.
The limitation of each platform is precisely the strength of the other: CARLA
lacks aerial agents, while AirSim lacks realistic ground traffic and pedestrian
interactions.
Meanwhile, AirSim's upstream development has been archived by its original
maintainers, leaving a widely adopted flight simulation stack without an active
evolution path.
Other simulators across driving, UAV, and embodied AI domains similarly remain
confined to a single agent modality (see Section~\ref{sec:related} for a
comprehensive survey).
As a result, emerging workloads that span both air and ground
domains---air-ground cooperation, cross-domain embodied navigation, joint
multi-modal data collection, and cooperative reinforcement learning---lack a
shared simulation foundation.

\paragraph{Why not bridge-based co-simulation?}
A common workaround connects heterogeneous simulators through bridge-based
co-simulation, typically via ROS\,2~\cite{ros2} or custom message-passing
interfaces.
While functionally viable, such approaches introduce inter-process
synchronization complexity, communication overhead, and duplicated rendering
pipelines.
More critically, independent simulation processes cannot guarantee strict
spatial-temporal consistency across sensor streams---a requirement for
perception, learning, and evaluation in embodied intelligence systems.
Fig.~\ref{fig:ipc_overhead} quantifies the per-frame inter-process overhead
contrast between bridge-based co-simulation and the single-process design
adopted by CARLA-Air.

\input{figures/ipc_plot}

\paragraph{CARLA-Air: a unified infrastructure for air-ground embodied intelligence.}
We present \textbf{CARLA-Air}, an open-source platform that integrates CARLA
and AirSim within a single Unreal Engine process, purpose-built as a practical
simulation foundation for air-ground embodied intelligence research.
By inheriting and extending AirSim's aerial simulation capabilities within a
modern, actively maintained infrastructure, CARLA-Air also provides a
sustainable evolution path for the large body of existing AirSim-based research.
Key capabilities of the platform include:

\begin{enumerate}[leftmargin=*, label=(\roman*)]

\item \textbf{Single-process air-ground integration.}
CARLA-Air resolves a fundamental engine-level conflict---UE4 permits only one
active game mode per world---through a composition-based design that inherits
all ground simulation subsystems from CARLA while spawning AirSim's aerial
flight actor as a regular world entity.
This yields a shared physics tick, a shared rendering pipeline, and strict
spatial-temporal consistency across all sensor viewpoints.

\item \textbf{Full API compatibility and zero-modification code migration.}
Both CARLA and AirSim native Python APIs and ROS\,2 interfaces are fully
preserved, allowing existing research codebases to run on CARLA-Air without
modification.

\item \textbf{Photorealistic, physically coherent simulation world.}
The platform delivers rich urban and natural environments populated with
rule-compliant traffic flow, socially-aware pedestrians, and aerodynamically
consistent multirotor dynamics, with synchronized capture of up to 18 sensor
modalities across all aerial and ground platforms at each simulation tick.

\item \textbf{Extensible asset pipeline.}
Researchers can import custom robot platforms, UAV configurations, vehicles,
and environment maps into the shared simulation world, enabling flexible
construction of diverse air-ground interaction scenarios.

\end{enumerate}

\noindent Building on these capabilities, CARLA-Air provides out-of-the-box
support for representative air-ground embodied intelligence workloads across
four research directions:

\begin{enumerate}[leftmargin=*, label=(\alph*)]

\item \textbf{Air-ground cooperation}---heterogeneous aerial and ground agents
coordinate within a shared environment for tasks such as cooperative
surveillance, escort, and search-and-rescue.

\item \textbf{Embodied navigation and vision-language action}---agents navigate
and act grounded in visual and linguistic input, leveraging both aerial
overview and ground-level detail.

\item \textbf{Multi-modal perception and dataset construction}---synchronized
aerial-ground sensor streams are collected at scale to build paired datasets
for cross-view perception, 3D reconstruction, and scene understanding.

\item \textbf{Reinforcement-learning-based policy training}---agents learn
cooperative or individual policies through closed-loop interaction in
physically consistent air-ground environments.

\end{enumerate}

\noindent As a lightweight and practical infrastructure, CARLA-Air lowers the
barrier for developing and evaluating air-ground embodied intelligence systems,
and provides a unified simulation foundation for emerging applications in
low-altitude robotics, cross-domain autonomy, and large-scale embodied AI
research.

